\documentclass[aspectratio=169,10pt]{beamer}
\usepackage{booktabs}
\usepackage{graphicx}
\setbeamertemplate{navigation symbols}{}
\setbeamercolor{frametitle}{fg=blue!60!black}
\setbeamertemplate{footline}[frame number]
\title{DCB+Gaussian versus sum of Gaussians}
\subtitle{Signal-model checks and smearing response: provisional comparison}
\date{19 October 2026}
\begin{document}
\begin{frame}{How should we choose the signal PDF?}
\textbf{Choose from agreement with signal simulation and fit stability.}
\begin{itemize}
\item Compare the same process, category and era at the generated MC mass, $m_H=125$ GeV; inspect the peak, tails, residuals, convergence and boundaries.
\item A goodness-of-fit probability needs a calibrated statistic, with weighted MC uncertainties and fitted parameters accounted for. Our summed diagnostic $\chi^2$ values are not calibrated model-selection tests.
\item Expected sensitivity is a consequence to compare, not the selection criterion.
\item If both families adequately describe MC, DCB+Gaussian is a reasonable simpler workflow: it avoids selecting the number of Gaussians with the SoG F-test. This does not imply fewer fitted coefficients in every case.
\end{itemize}
\small Existing all-era aggregate diagnostic: lower $\chi^2$ for DCB in 7 categories and SoG in 9; no universal winner. Component failures can be hidden by a good-looking sum.
\vfill\tiny Source: Signal/packaged\_all\_years\_20261014/README\_MH125.md and aggregate\_comparison\_MH125.csv.
\end{frame}

% BEGIN SHAPE EXAMPLES
\begin{frame}{Visible shape difference: tH\_had\_4, 2023postBPix}
\begin{columns}[T]
\column{0.49\textwidth}\centering\textbf{DCB+Gaussian}\\
\includegraphics[height=0.66\textheight]{../runplotter_individual_years_20261014/DCB/2023postBPix/Plots/smodel_tH_had_4_2023postBPix.png}
\column{0.49\textwidth}\centering\textbf{Sum of Gaussians}\\
\includegraphics[height=0.66\textheight]{../runplotter_individual_years_20261014/SoG/2023postBPix/Plots/smodel_tH_had_4_2023postBPix.png}
\end{columns}
\small SoG shows a shoulder near 128 GeV; DCB has a smoother falling tail. This motivates component-level checks; sparse MC does not establish that every shoulder is unphysical.
\vfill\tiny Matched era/category; sum of processes, $M_H=m_H^{MC}=125$ GeV, nominal nuisances, 0.1\% display pruning. The smoothness comparison is not a calibrated goodness-of-fit verdict. Source: runplotter\_individual\_years\_20261014.
\end{frame}
\begin{frame}{Visible shape difference: tH\_had\_3, 2023postBPix}
\begin{columns}[T]
\column{0.49\textwidth}\centering\textbf{DCB+Gaussian}\\
\includegraphics[height=0.66\textheight]{../runplotter_individual_years_20261014/DCB/2023postBPix/Plots/smodel_tH_had_3_2023postBPix.png}
\column{0.49\textwidth}\centering\textbf{Sum of Gaussians}\\
\includegraphics[height=0.66\textheight]{../runplotter_individual_years_20261014/SoG/2023postBPix/Plots/smodel_tH_had_3_2023postBPix.png}
\end{columns}
\small SoG has a flatter peak and a low-side shoulder; DCB has a smoother dominant peak. Similar effective widths (2.06 versus 2.05 GeV) do not imply identical local shapes.
\vfill\tiny Matched era/category and plotting prescription as preceding slide. These are process sums, not individual fit curves; their appearance alone does not identify a failed donor or prove model adequacy. Source: runplotter\_individual\_years\_20261014.
\end{frame}
\begin{frame}{A documented poor SoG solution: inspect residuals and fit status}
\begin{columns}
\column{0.66\textwidth}
\includegraphics[width=\linewidth,height=0.75\textheight,keepaspectratio]{../shape_review_20261014/2022preEE_bbh_in_tH_lep_2.png}
\column{0.32\textwidth}\small
2022preEE, recipient \texttt{bbh\_in / tH\_lep\_2}; both borrow \texttt{tth\_in / tH\_lep\_2}.
\begin{itemize}
\item SoG peak shifts relative to donor MC; structured residuals reach about 5--7 error units.
\item Logged $\chi^2/\mathrm{ndof}$: DCB 1.050; SoG 14.317.
\item DCB optimizer converged, with two bound flags; SoG optimizer failed.
\end{itemize}
\textbf{Conclusion:} the returned SoG solution needs correction; this does not show that a converged SoG cannot fit this shape.
\end{columns}
\vfill\tiny Normalized donor-shape comparison at $M_H=125$ GeV; not a recipient-yield comparison. Reduced $\chi^2$ is diagnostic, not a calibrated p-value. Source: shape\_review\_20261014/selected\_shapes.json and README.md.
\end{frame}
\begin{frame}{Keep the comparison balanced: a case where SoG is better}
\begin{columns}
\column{0.66\textwidth}
\includegraphics[width=\linewidth,height=0.75\textheight,keepaspectratio]{../shape_review_20261014/2023postBPix_tHW_in_ttH_lep_1.png}
\column{0.32\textwidth}\small
2023postBPix, \texttt{tHW\_in / ttH\_lep\_1}.
\begin{itemize}
\item Both optimizers converged.
\item SoG follows the MC peak and residuals more closely in this example.
\item Logged $\chi^2/\mathrm{ndof}$: DCB 2.379; SoG 0.655.
\item DCB has two bound flags; SoG has none.
\end{itemize}
\textbf{Selection principle:} review fit quality, convergence and tails across matched components; do not select by sensitivity or visual smoothness alone.
\end{columns}
\vfill\tiny Selected illustrative cases are not an unbiased survey of model quality. Source: shape\_review\_20261014/selected\_shapes.json. Existing examples do not establish a universal winning family.
\end{frame}
% END SHAPE EXAMPLES

\begin{frame}{How does the smearing nuisance change the signal shape?}
\begin{columns}
\column{0.64\textwidth}
\includegraphics[width=\linewidth,height=0.73\textheight,keepaspectratio]{../smearing_pdf_scan_20261015/2023preBPix_peak.png}
\column{0.34\textwidth}
\small Matched example: \texttt{tth\_in}, \texttt{ttH\_had\_1}, 2023preBPix; $M_H=125.09$ GeV.
\begin{itemize}
\item Both nominal optimizers converged in this example.
\item At $\theta=\pm1$, widths change by $\pm4.36\%$ in both families; there is also a small mean response ($\sim0.037$ GeV).
\item Tail coefficients and mixture fractions are not refitted or directly varied by this nuisance.
\item PDF response is distinct from the profiled signal-strength impact.
\end{itemize}
\end{columns}
\tiny A SoG with distinct component means need not be symmetric. DCB tails allow left/right differences; neither family guarantees symmetric or asymmetric impacts. Source: smearing\_pdf\_scan\_20261015/README.md; tools/finalModel.py, buildSigma/buildMean.
\end{frame}
\begin{frame}{Historical DCB change: a separate question}
\small
\begin{center}
\begin{tabular}{lrrr}\toprule
Result & $\Delta r_{ttH}(\theta_-)$ & $\Delta r_{ttH}(\theta_+)$ & Rank\\\midrule
Earlier DCB & -0.0180 & +0.0313 & 6 \\
Current DCB & -0.0192 & +0.0183 & 4 \\
\bottomrule\end{tabular}
\end{center}
\begin{itemize}
\item The current DCB impact is nearly symmetric and fourth among the 91 collected constrained nuisances. The earlier result is more asymmetric.
\item Ranking depends on all other impacts too. These numbers do not establish why the result changed; earlier and current fit setups are not assumed identical.
\item Run fresh SoG Combine fits with matched mass, POIs, data-generation prescription, constraints, background envelope, pruning and numerical settings. Check Asimov closure and fit quality before comparing impacts.
\item Compare absolute up/down shifts, nuisance pulls/constraints and the PDF response, not just the ranking. For expected sensitivity, a shared reference Asimov dataset is useful to separate shape-model changes from changes in generated data.
\end{itemize}
\textbf{Update:} current SoG results follow. A controlled explanation of the historical DCB change remains pending.
\vfill\tiny Current DCB caveat: initial snapshot status 1, covariance quality 2; one impact log also needs stability review. Source: Combine/DCB/impacts/collected\_20261017/impacts.json and README.md; earlier raw JSON in Combine/15OCT26/DCB\_0p001/.
\end{frame}

% BEGIN CURRENT COMBINE COMPARISON
\begin{frame}{Current impacts: smearing is fourth for ttH in both families}
\begin{center}\small
\begin{tabular}{llrrr}\toprule
POI & Model & Lower-endpoint shift & Upper-endpoint shift & Rank\\\midrule
r\_ttH & DCB & -0.0192 & +0.0183 & 4 \\
r\_ttH & SoG & -0.0178 & +0.0172 & 4 \\
r\_tHq & DCB & -0.0669 & +0.0751 & 12 \\
r\_tHq & SoG & -0.0551 & +0.0666 & 13 \\
\bottomrule\end{tabular}\end{center}
\begin{center}\includegraphics[width=0.87\linewidth,height=0.40\textheight,keepaspectratio]{smearing_impact_comparison.png}\end{center}
\small Both ttH impacts are nearly symmetric; DCB is modestly larger. This result does not support a general rule that SoG impacts must be symmetric and DCB impacts asymmetric.
\vfill\tiny All 91 constrained nuisances collected in each model, other strength profiled. Endpoints use fitted nuisance intervals, not the bare conditional shape scan. Original JSON values retained; no nuisance removal or recentering.
\end{frame}
\begin{frame}{Expected profile scans: very similar over the relevant intervals}
\begin{center}\includegraphics[width=\linewidth,height=0.75\textheight,keepaspectratio]{combine_scan_comparison.png}\end{center}
\small Solid: total uncertainties. Dashed: stat only. Raw likelihoods; invalid DCB tail points omitted without bridging gaps. SoG refined scans have no invalid grid points.
\vfill\tiny Displayed ranges overlap; scan grids differ. Each family uses its own saved SM-injected Asimov. Snapshot status/covariance warnings remain; these are provisional expected comparisons, not a common-dataset closure test.
\end{frame}
\begin{frame}{Numerical comparison and what it establishes}
\begin{center}\small
\begin{tabular}{lcccc}\toprule
Model & ttH total 68\% interval & tHq total 68\% interval & ttH $Z$ & tHq median limit\\\midrule
DCB & [0.615, 1.535] & [-3.507, 6.535] & 2.827 & 11.445 \\
SoG & [0.615, 1.534] & [-3.504, 6.531] & 2.831 & 11.445 \\
\bottomrule\end{tabular}\end{center}
\begin{itemize}
\item Expected POI intervals, significances and the reported tHq upper limit are nearly unchanged. The limit is a dimensionless strength, at 95\% CLs.
\item Smearing is fourth for ttH in both models. Its smaller SoG impact does not by itself select the preferred signal PDF.
\item Base PDF choice on signal-MC agreement and stability. The shown SoG shape failures and DCB counterexamples need component-level review.
\item Remaining qualifications: DCB snapshot status 1/covariance 2; SoG status 3/covariance 3. Both best fits are near $(1.042,1.020)$, rather than exact closure at $(1,1)$.
\item These are matched-workflow, separate-Asimov results. A shared reference dataset is still needed to isolate the fitted-family effect; the cause of the historical DCB impact change is not established.
\end{itemize}
\vfill\tiny Sources: Combine/DCB and Combine/SoG saved scan summaries, significance/limit results and collected impacts; reproducible table inputs in combine\_comparison.json.
\end{frame}
% END CURRENT COMBINE COMPARISON

\end{document}
